\chapter{Introduction}
%---------------------------------------------------------------------
\section{Endosymbiosis and its importance in eukaryotic evolution}
Endosymbiosis can be defined as a relationship between two organisms in which one of them, the \gls{ES}, resides in the other, the host. It can provide physiological advantages to the host, e.g., upscaling of cellular energy or enabling hosts to survive on nutritionally imbalanced diets \parencite{bennett_endosymbioses_2024, sorensen_protein_2025}. Further, it has shaped the evolutionary trajectory of eukaryotes and continues to do so, resulting in a large diversity of endosymbiotic relationships \parencite{douglas_symbiosis_2014, archibald_endosymbiosis_2015, sorensen_protein_2025, bennett_endosymbioses_2024}. This underlines its importance in the history of life on Earth. 
\par It is hypothesized that an ancestor to the Asgard archaea, in its path to become the eukaryotic cell lineage, acquired and integrated a free-living bacterium as an \gls{ES} \parencite{spang_complex_2015, zaremba-niedzwiedzka_asgard_2017}. Over an evolutionary time frame, the acquisition of the \gls{ES} is accompanied by the following phenomena: (i) a metabolic exchange and flux control, whereby the \gls{ES} becomes metabolically streamlined \parencite{bennett_endosymbioses_2024, nowack_genomics-informed_2018}, (ii) host control over the \gls{ES} ensuring synchronous cell division \parencite{sorensen_protein_2025, morales_host-symbiont_2023}, (iii) genome reduction of the \gls{ES} and \gls{EGT} events resulting in proteins originally from the \gls{ES} becoming encoded by the nucleus \parencite{nowack_genomics-informed_2018, timmis_endosymbiotic_2004}, and lastly (iv) protein targeting and import mechanisms for the nucleus encoded proteins \parencite{sorensen_protein_2025, nowack_genomics-informed_2018}. This process can be summarized as organellogenesis, though the line between \gls{ES} and organelle is fluid, e.g., spheroid bodies inside diatomes \parencite{nowack_endosymbiotic_2010, archibald_endosymbiosis_2015}.
\par To date, only four primary endosymbiosis events involving organellogenesis have been described in the literature. The two canonical organelles, the mitochondrion and the chloroplast, developed $\sim$1.8 billion years ago from an \textgreek{α}-proteobacterium and $\sim$1.5 billion years ago from a \textgreek{β}-cyanobacterium, respectively \parencite{parfrey_estimating_2011, gray_mitochondrial_1999, yoon_molecular_2004}. The other two incidents of organellogenesis comprimise: (i) the chromatophore of \emph{Paulinella chromatophora} which evolved from an \textgreek{α}-cyanobacterium around 90 -- 140 million years ago \parencite{nowack_paulinella_2014}, and (ii) the potential candidate for the nitroplast, the \gls{ucyna}.
%---------------------------------------------------------------------
%%%%%%%%%%%%%%%%%%%%%%%%%%%%%%%%%%%%%%%%%%%%%%%%%%%%%%%%%%
%---------------------------------------------------------------------
\section{The nitroplast: A recently recognized organelle}
\gls{ucyna} is the \gls{ES} of \glsxtrlong{Bbig},  a species of phytoplankton belonging to the group of the haptophytes \parencite{hagino_discovery_2013}. Notably, \gls{ucyna} has retained a thin peptidoglycan wall and is surrounded by three membranes, one from the host and two bacterial membranes \parencite{hagino_discovery_2013}. Along with \gls{ucyna}, \glsxtrshort{Bbig} harbors mitochondria and two secondary plastids.
Moreover, \glsxtrshort{Bbig} has two distinct ecotypes: the coastal-form that produces calcium-carbonate shells and the open-ocean-form that has the haptonema for motility \parencite{hagino_discovery_2013, thompson_genetic_2014}. Interestingly, the two ecotypes also differ in the size of the \gls{ES} as well as in the number of endosymbionts present per cell, and it has been suggested that these ecotypes coevolved with their respective \gls{ES} for at least $\sim$91 million years, diverging from each other and adapting to their environments \parencite{cornejo-castillo_cyanobacterial_2016, thompson_genetic_2014}. 
Furthermore, \gls{ucyna} has been classified as the first N\textsubscript{2}-fixing organelle, the nitroplast, for the following reasons. It has a reduced genome of 1.4 Mbp lacking genes for the tricarboxylic acid cycle, the carbon fixation, as well as the photosynthesis system II, but has retained the entire set of \textit{nif} genes,  indicating metabolic streamlining for N\textsubscript{2}-fixation \parencite{tripp_metabolic_2010, zehr_globally_2008, cornejo-castillo_cyanobacterial_2016}. Though no \gls{EGT} has been observed to date \parencite{suzuki_unstable_2021}. In addition to metabolic streamlining, \gls{ucyna} is integrated into the cellular architecture as it displays an organelle size relationship and divides synchronously with the host \parencite{coale_nitrogen-fixing_2024, cornejo-castillo_metabolic_2024}. Most interestingly, large-scale protein targeting has been described for \gls{ucyna}.

%---------------------------------------------------------------------
%%%%%%%%%%%%%%%%%%%%%%%%%%%%%%%%%%%%%%%%%%%%%%%%%%%%%%%%%%
%---------------------------------------------------------------------
\section{Protein targeting and protein import}
\cite{coale_nitrogen-fixing_2024} demonstrated that some proteins encoded by \glsxtrshort{Bbig} have been enriched in \gls{ucyna}, indicating protein import. Through motif finding tools, the team has discovered that these proteins share an extension of around 120 amino acids at the C-terminal region, which has not been detected in the proteome analysis, suggesting that it is cleaved off once the protein is inside the nitroplast. This C-terminal sequence has been declared as the \gls{utp}. The mechanism behind its specificity in targeting and import into the nitroplast is yet unknown.
\par Transit peptides have been described for the other aforementioned organelles. \Gls{mtps} and \gls{ctps} share some commonalities. They are highly diverse in their sequence but display conserved physicochemical properties, like containing one or more α-helices with amphiphathic character \parencite{schleiff_common_2011}. The transit peptides have been described to consist of two domains: an N-terminal region conveying specificity of targeting and a C-terminal region that facilitates translocation \parencite{lee_molecular_2019}. Proteins tagged by these two transit peptides are transported in an unfolded state with the help of chaperons to their organelle, where they are imported via translocase complexes of the outer and inner membranes into the organelle, followed by cleavage of the transit peptide \parencite{schleiff_common_2011}. Organelle targeting specificity is facilitated by positive charges in the \gls{mtps} \parencite{kunze_similarity_2015}.
\par In contrast, \gls{crtps} are more distinct from \gls{mtps} and \gls{ctps}. The \gls{crtps} are much larger, around 200 amino acids, compared to the 15~--~80 amino acids for the canonical transit peptides \parencite{garg_role_2016, von_heinje_chloroplast_1991, sorensen_protein_2025}. They are also bipartite. The crTP\textsubscript{part 1} consists of a transmembrane helix and the processing region 1, which contains the signals for the secretory pathway and is cleaved off during import \parencite{oberleitner_bipartite_2022, klimenko_paulinella_2025}. The crTP\textsubscript{part 2} is a γ-glutamyl cyclotransferase-like domain which folds after cleavage of part 1 and most likely mediates the passage of the protein through the peptidoglycan cell wall of the chromatophore  \parencite{klimenko_paulinella_2025, oberleitner_bipartite_2022}.
%---------------------------------------------------------------------
%%%%%%%%%%%%%%%%%%%%%%%%%%%%%%%%%%%%%%%%%%%%%%%%%%%%%%%%%%
%---------------------------------------------------------------------
\section{Heterologous expression}
%---------------------------------------------------------------------
\subsection{\textit{Angomonas deanei} and its endosymbiont}
\glsxtrshort{Adea} is a monoxenous parasite of insects that belongs to the the subfamily of the Strigomonadinae within the family of Trypanosomatidae \parencite{votypka_kentomonas_2014}. Like the other members in its subfamily it harbors a β-proteobacterium, \glsxtrlong{Kineto}, which originated from a single endosymbiosis event around 40~--~120 million years ago \parencite{du_monophyletic_1994, teixeira_phylogenetic_2011}. The coevolution of \glsxtrshort{Adea} and \glsxtrshort{Kineto} led to a reduction in the genome size of the \gls{ES} to 0.8 Mbp with one \gls{EGT} event (ornithine cyclodeaminase) that had occurred \parencite{alves_endosymbiosis_2013, alves_genome_2013}. This process was accompanied by the integration of the \gls{ES} in the metabolism of \glsxtrshort{Adea}, e.g., the complementation of the heme biosynthesis \parencite{klein_biosynthesis_2013}. Moreover, the cell division of \glsxtrshort{Kineto} and \glsxtrshort{Adea} occurs synchronously and is under the control of the host, such that only one \gls{ES} is present in every cell \parencite{motta_bacterium_2010, maurya_nucleus-encoded_2025}. Such control is mediated by proteins encoded by \glsxtrshort{Adea} that localize to the \gls{ES}; such a protein is referred to as an \gls{etp} and to this date, seven \gls{etp}s have been identified  \parencite{morales_development_2016, morales_host-symbiont_2023}. However, these \gls{etp}s do not possess a transit peptide or share many similarities. Interestingly, \gls{etp}3 and \gls{etp}8 display localization inside the \gls{ES}, suggesting protein import, yet it remains unknown how protein import occurs \parencite{morales_host-symbiont_2023}. Furthermore, the genomes of \glsxtrshort{Adea} and \glsxtrshort{Kineto} as well as the transcriptome of \glsxtrshort{Adea} are readily available \parencite{morales_development_2016, motta_predicting_2013,alves_genome_2013}. In addition to that,  molecular tools have been devised for the integration of heterologous genes inside the genome of \glsxtrshort{Adea} \parencite{morales_development_2016}. Thus, \glsxtrshort{Adea} and its \gls{ES} represent a useful model system for investigating protein targeting, especially involving the expression of heterologous transit peptides. 
%---------------------------------------------------------------------
%---------------------------------------------------------------------
\subsection{\textit{Escherichia coli} and purifications}
\glsxtrshort{Eco} is used as a platform for heterologous expression because of its rapid growth to high cell densities, high protein production at relatively low maintenance and costs, as well as the existence of high efficiency protocols for transformation \parencite{sezonov_escherichia_2007, shiloach_growing_2005, rosano_recombinant_2014}. Moreover, the most common expression strain BL21, which other engineered strains are based on, contains on its chromosome the \gls{t7rnap} under the control of a \textit{lac} promoter allowing for the inducible expression of plasmids containing the protein of interest under the control of a T7 promoter \parencite{rosano_recombinant_2014, terpe_overview_2006, studier_use_1986}.

Here, the goal of heterologous expression in \glsxtrshort{Eco} is to produce soluble recombinant \gls{utp}s containing an affinity tag at high concentrations, such that they can be purified for structural analysis. A common method for the purification of protein is \gls{IMAC}, which uses metals like nickel to bind a \gls{6xhis}-tag attached to the protein that can then be displaced by imidazole to elute the protein of interest \parencite{porath_immobilized_1992}. Once enough purified protein has been obtained, structural analysis can be done to gain information about the protein folding and, therefore, its function.
%---------------------------------------------------------------------
%%%%%%%%%%%%%%%%%%%%%%%%%%%%%%%%%%%%%%%%%%%%%%%%%%%%%%%%%%
%---------------------------------------------------------------------
\section{Aim of this study}
Since the establishment of the nitroplast as well as the endosymbiotic system of \glsxtrshort{Adea} with its \gls{ES} are relatively young compared to the canonical organelles they can serve as models for uncovering the evolution of protein targeting. In this work, three \gls{utp}s (\glsxtrshort{utp_B}, \glsxtrshort{utp_E}, \glsxtrshort{utp_C}) have been selected to be investigated in the different workflows.

\par The first part of this work aims to understand the impact of recombinant nitroplast targeting signals on different endosymbiotic systems, here \glsxtrshort{Adea}. For that purpose, three different \gls{utp}s will be tagged with an \gls{egfp}, the constructs will then be integrated and expressed in \glsxtrshort{Adea}. Considering that protein targeting to the \gls{ES} in \glsxtrshort{Adea} lacks specific signals, it can be hypothesized that these \gls{egfp}-tagged \gls{utp}s  will most likely remain in the cytosole of \gls{Adea}. A Western blot analysis and in-gel fluorescence assay have been performed to verify expression of correctly folded \gls{egfp}-\gls{utp} constructs in \glsxtrshort{Adea}. Additionally, epifluorescence and confocal fluorescence microscopy have been done to localize these fusion proteins within \glsxtrshort{Adea} and see whether or not they colocalize with the \gls{ES}-marker, \glsxtrlong{mSc}-\gls{etp}1. Evaluating the nitroplast transit peptides based on the effects of heterologous expression and localization within \glsxtrshort{Adea} would provide insights into the possible similarities or even uncover universal principles guiding the different targeting mechanisms.

\par The second part of this work intends to test the expression and purification of these three \gls{utp}s in \glsxtrshort{Eco} to obtain purified proteins for structural analysis. A \gls{6xhis} tag with \gls{sumo} as well as a \gls{tev site} has been cloned to the N-terminus of the \gls{utp}s, and \glsxtrshort{Eco} LOBSTR has been transformed with these constructs. Here, expression, solubility and purification tests were conducted with \gls{utp} fusion proteins. It can be predicted that heterologous expression as well as purification should be possible. Expression and solubility tests have been utilized to attain an expression condition yielding the most soluble fusion proteins. When a condition has been chosen, a purification test using \gls{IMAC} was done to examine the possibility of purifying the \gls{utp}s containing the \gls{6xhis} tag. Based on the results of the test, larger-scale expression and purification using \gls{fplc} can be done to obtain enough purified protein for structural analysis, thus laying the groundwork for uncovering the structures of the nitroplast targeting signals.

\par Additionally, bioinformatic analyses have been utilized to gather information on the structure and properties of the selected \gls{utp}s. A sequence alignment with Clustal Ω has been done. Structures have also been predicted using Alphafold 3. Further analysis tools, like HeliQuest and DeepLoc, have also been used. Using these tools will likely facilitate our understanding of the structure, function, and targeting mechanism of the \gls{utp}s. 
